\documentclass[11pt,a4paper]{article}

%========================
% PACKAGES
%========================
\usepackage[utf8]{inputenc}
\usepackage[T1]{fontenc}
\usepackage[french]{babel}
\usepackage[a4paper,margin=2cm]{geometry}
\usepackage{amsmath,amssymb}
\usepackage{lmodern}
\usepackage{fancyhdr}
\usepackage{lastpage}
\usepackage{enumitem}

%========================
% MISE EN PAGE
%========================
\pagestyle{fancy}
\fancyhf{}
\fancyhead[L]{Mathématiques}
\fancyhead[C]{Devoir surveillé -- Calcul numérique}
\fancyhead[R]{Page \thepage/\pageref{LastPage}}

%========================
\begin{document}

\begin{center}
\Large \textbf{DEVOIR SURVEILLÉ}\\
\large Calcul numérique
\end{center}

\vspace{0.5cm}

\textbf{Durée :} 1h30 \hfill \textbf{Calculatrice autorisée}\\
\textbf{Total : 28 points}

\textbf{Consigne :} Justifier tous les calculs.

\vspace{0.5cm}

%========================
\section*{Exercice 1 -- Problème contextualisé (type brevet)}

Un parc de loisirs propose des entrées à la journée.

\begin{itemize}
\item Le prix d’une entrée adulte est de 24€.
\item Le prix d’une entrée enfant est de 16€.
\item Une famille achète 5 billets pour un total de 104€.
\end{itemize}

\begin{enumerate}
\item Montrer que le nombre d’adultes et d’enfants vérifie le système : \hfill (2 pts)
\[
\begin{cases}
a + e = 5 \\
24a + 16e = 104
\end{cases}
\]

\item Résoudre ce système et déterminer le nombre d’adultes et d’enfants. \hfill (3 pts)

\item Le parc propose ensuite une réduction de 20\% sur le total.
\begin{enumerate}
\item Calculer le montant de la réduction. \hfill (0.5 pt)
\item Calculer le prix final payé. \hfill (1 pt)
\end{enumerate}

\item Une autre famille paie 96€ après réduction.
\begin{enumerate}
\item Retrouver le prix avant réduction. \hfill (1 pt)
\item Combien de billets ont-ils achetés s’ils n’ont pris que des billets enfants ? \hfill (0.5 pt)
\end{enumerate}

\end{enumerate}

\vspace{0.5cm}

%========================
\section*{Exercice 2 -- Calculs sur les fractions}

Dans cet exercice, on cherche à comprendre les \textbf{règles de calcul sur les fractions}.

\subsection*{Partie A -- Calculs}

\begin{enumerate}
\item Calculer :
\[
A = \frac{3}{4} + \frac{5}{6}
\] \hfill (0.5 pt)

\item Calculer :
\[
B = \frac{7}{5} - \frac{2}{3}
\] \hfill (0.5 pt)

\item Calculer :
\[
C = \frac{3}{4} \times \frac{8}{9}
\] \hfill (0.5 pt)

\item Calculer :
\[
D = \frac{5}{6} \div \frac{2}{3}
\] \hfill (0.5 pt)

\end{enumerate}

\subsection*{Partie B -- Erreurs classiques}

On considère les affirmations suivantes :

\begin{enumerate}
\item $\displaystyle \frac{2}{3} + \frac{1}{3} = \frac{3}{6}$ \hfill (1 pt)

\item $\displaystyle \frac{3}{4} + \frac{1}{2} = \frac{4}{6}$ \hfill (1 pt)

\item $\displaystyle \frac{2}{5} \times \frac{3}{7} = \frac{6}{35}$ \hfill (1 pt)

\item $\displaystyle \frac{4}{9} \div \frac{2}{3} = \frac{4}{9} \times \frac{3}{2}$ \hfill (1 pt)
\end{enumerate}

\begin{itemize}
\item Dire si chaque affirmation est vraie ou fausse.
\item Justifier en expliquant \textbf{la règle correcte}.
\end{itemize}

\subsection*{Partie C -- Calcul complexe}

Calculer :
\[
E = \frac{\frac{3}{4} + \frac{1}{2}}{\frac{5}{6} - \frac{1}{3}}
\] \hfill (2 pts)

\vspace{0.5cm}

%========================
\section*{Exercice 3 -- Problème choisi (raisonnement)}

On considère un nombre $x$.

On définit l’expression :
\[
F(x) = \frac{2x + 3}{x - 1}
\]

\begin{enumerate}
\item Calculer $F(2)$ et $F(3)$. \hfill (1 pt)

\item Peut-on calculer $F(1)$ ? Justifier. \hfill (1 pt)

\item On suppose que :
\[
F(x) = 2 + \frac{5}{x - 1}
\]

\begin{enumerate}
\item Vérifier cette égalité. \hfill (2 pts)
\item En déduire une méthode rapide pour calculer $F(101)$. \hfill (0.5 pt)
\end{enumerate}

\item Trouver une valeur de $x$ telle que :
\[
F(x) = 2
\] \hfill (0.5 pt)

\end{enumerate}

\vspace{1cm}

%========================
\section*{Exercice 4 -- Redémonstration des règles (4 points)}

\textbf{Objectif : reconstruire les règles de calcul sur les fractions.}

\textbf{Barème :}
\begin{itemize}
\item Rigueur des démonstrations : 3 pts
\item Qualité des explications : 1 pt
\end{itemize}

Soient $\frac{a}{b}$ et $\frac{c}{d}$ deux fractions avec $b \neq 0$ et $d \neq 0$.

\subsection*{Partie A -- Addition de fractions}

On cherche à démontrer que :
\[
\frac{a}{b} + \frac{c}{d} = \frac{ad + bc}{bd}
\]

\begin{enumerate}
\item Pourquoi ne peut-on pas additionner directement les dénominateurs ? \hfill (0.5 pt)

\item Quel est le dénominateur commun le plus simple entre $b$ et $d$ ? \hfill (0.5 pt)

\item Compléter :
\[
\frac{a}{b} = \frac{a \times \dots}{b \times \dots}
\quad \text{et} \quad
\frac{c}{d} = \frac{c \times \dots}{d \times \dots}
\] \hfill (0.5 pt)

\item Écrire $\frac{a}{b}$ avec le dénominateur $bd$. \hfill (0.5 pt)

\item Écrire $\frac{c}{d}$ avec le dénominateur $bd$. \hfill (0.5 pt)

\item Additionner les deux fractions obtenues. \hfill (0.5 pt)

\item Simplifier l’expression et conclure. \hfill (0.5 pt)

\end{enumerate}

\vspace{0.3cm}

\subsection*{Partie B -- Multiplication de fractions}

On cherche à montrer que :
\[
\frac{a}{b} \times \frac{c}{d} = \frac{ac}{bd}
\]

\begin{enumerate}
\item Écrire $\frac{a}{b}$ comme $a \times \frac{1}{b}$. \hfill (0.5 pt)

\item Écrire $\frac{c}{d}$ comme $c \times \frac{1}{d}$. \hfill (0.5 pt)

\item Regrouper tous les facteurs. \hfill (0.5 pt)

\item Simplifier pour obtenir le résultat. \hfill (0.5 pt)

\end{enumerate}

\vspace{0.3cm}

\subsection*{Partie C -- Division de fractions}

On cherche à montrer que :
\[
\frac{a}{b} \div \frac{c}{d} = \frac{a}{b} \times \frac{d}{c}
\]

\begin{enumerate}
\item Que signifie ``diviser par un nombre'' en mathématiques ? \hfill (0.5 pt)

\item Quel est l’inverse de $\frac{c}{d}$ ? \hfill (0.5 pt)

\item Réécrire la division comme une multiplication. \hfill (0.5 pt)

\item Conclure. \hfill (0.5 pt)

\end{enumerate}

\vspace{0.3cm}

\subsection*{Partie D -- Compréhension}

\begin{enumerate}
\item Expliquer avec tes mots pourquoi on ne peut pas faire :
\[
\frac{a}{b} + \frac{c}{d} = \frac{a+c}{b+d}
\] \hfill (0.5 pt)

\item Donner un contre-exemple numérique simple. \hfill (0.5 pt)

\end{enumerate}

\vspace{1cm}

\begin{center}
\textbf{Fin du sujet}
\end{center}

\end{document}